

% However, current polarization observations are only available at a single wavelength. Multi-wavelength polarization observations are needed to determine the dominant polarization mechanism and place constraints on the dust grain size. 


Although previous studies of IRS 63 have revealed a structured disk and detected polarized emission at \bsixNS, the disk has not been explored using multi-wavelength polarization observations. These data are essential to disentangle the polarization morphologies with wavelength and determine the mechanisms polarizing the light. It can also provide a constraint on the dust grain sizes, providing an insight into the early stages of planet formation. 


The main goal of this thesis is to use multi-wavelength polarization observations of IRS 63 to constrain the dust grain size in the disk. To achieve this main goal, this thesis has the following objectives:
\begin{enumerate}
    \item Characterize and geometrically model the polarization morphologies of IRS 63 at \bfourNS, \bfiveNS, \bsixNS, and \bsevenNS. 
    \item Model the observed polarization using dust scattering models with different grain sizes, shapes and porosities. 
    \item Constrain the maximum dust grain size and compare to dust grain sizes of the cloud core and ISM to look for evidence of dust grain growth. 
\end{enumerate}
